\documentclass{studentnotes}
\usepackage{siunitx}
\usepackage{multicol}

\setnotetitle{Resonance, Transients, and Wave Propagation}
\setnotecourse{Continuous Notes on Coupled SHM and Waves}
\setnoteauthor{Ot}
\setnotedate{\today}
\lhead{\textbf{Resonance and Waves}}
\chead{Continuous Study Note}
\rhead{\thepage}

\begin{document}

\makenotetitle

\tableofcontents
\newpage

\section{From One Oscillator to a Wave}

An isolated simple harmonic oscillator is the basic unit. A coupled system is
what happens when two or more such units can exchange energy. A wave is what
happens when many coupled units form a continuous medium, so that a disturbance
started in one region is passed on to neighbouring regions.

The common mathematical thread is a restoring effect plus inertia. The restoring
effect tries to bring the system back toward equilibrium; inertia prevents the
return from being instantaneous. This delay is what makes oscillation possible.
When oscillators are coupled, energy can move from one oscillator to another. In
a continuous medium, that transfer appears as propagation.

\begin{importantnote}
    The most important distinction in this note is between \textbf{natural
    behaviour} and \textbf{forced behaviour}. Natural behaviour depends on the
    system itself. Forced behaviour also depends on the driving source.
\end{importantnote}

\section{Coupled SHM Systems}

Consider two equal masses \(m\), each attached to fixed supports by springs of
constant \(k\), and coupled to each other by a spring of constant \(k_c\). Let
\(x_1\) and \(x_2\) be the small displacements of the masses from equilibrium.
The equations of motion can be written as

\[
    m\ddot{x}_1=-kx_1-k_c(x_1-x_2),
    \qquad
    m\ddot{x}_2=-kx_2-k_c(x_2-x_1).
\]

The coupled terms contain \(x_1-x_2\), because the coupling spring cares about
relative displacement. If the masses move together, the coupling spring is not
stretched. If they move oppositely, the coupling spring is stretched strongly.

\begin{namedformula}{Two identical coupled oscillators}
    \ddot{x}_1+\frac{k+k_c}{m}x_1-\frac{k_c}{m}x_2=0,
    \qquad
    \ddot{x}_2+\frac{k+k_c}{m}x_2-\frac{k_c}{m}x_1=0
\end{namedformula}

The best coordinates are not always \(x_1\) and \(x_2\). Instead, use normal
coordinates:

\begin{namedformula}{Normal coordinates}
    q_+=x_1+x_2,
    \qquad
    q_-=x_1-x_2
\end{namedformula}

The \(q_+\) mode describes the masses moving in phase. The \(q_-\) mode describes
the masses moving out of phase. Substituting these combinations into the
equations gives two independent SHM equations.

\begin{namedformula}{Normal-mode frequencies}
    \omega_+=\sqrt{\frac{k}{m}},
    \qquad
    \omega_-=\sqrt{\frac{k+2k_c}{m}}
\end{namedformula}

The in-phase mode has the lower frequency because the coupling spring is not
stretched. The out-of-phase mode has the higher frequency because the coupling
spring adds extra restoring force.

\subsection{Worked Example: Normal Modes}

Let \(m=\SI{0.50}{kg}\), \(k=\SI{18}{N.m^{-1}}\), and
\(k_c=\SI{2.0}{N.m^{-1}}\). Then

\[
    \omega_+=\sqrt{\frac{18}{0.50}}=\SI{6.0}{rad.s^{-1}}.
\]

\[
    \omega_-=\sqrt{\frac{18+2(2.0)}{0.50}}
             =\sqrt{44}
             =\SI{6.63}{rad.s^{-1}}.
\]

\[
    \boxed{\omega_+=\SI{6.0}{rad.s^{-1}},
    \qquad
    \omega_-=\SI{6.6}{rad.s^{-1}}}
\]

The two natural frequencies are close because the coupling spring is weak
compared with the support springs.

\subsection{Energy Exchange and Beats}

If only one oscillator is displaced at \(t=0\), the motion is usually not a pure
normal mode. It is a superposition of the in-phase and out-of-phase modes. The
two normal frequencies then interfere in time, producing a slow transfer of
energy back and forth between the oscillators.

When \(\omega_+\) and \(\omega_-\) are close, the energy exchange is slow. The
beat angular frequency is approximately the difference between the two normal
angular frequencies.

\begin{namedformula}{Beat rate in weakly coupled oscillators}
    \omega_{\text{beat}}=|\omega_- - \omega_+|
\end{namedformula}

In formula F2.4, \(\omega_{\text{beat}}\) is the angular frequency of the slow
envelope, or the rate at which energy appears to move back and forth between
the two oscillators. The symbols \(\omega_-\) and \(\omega_+\) are the angular
frequencies of the two normal modes. The vertical bars mean ``take the positive
difference'', because a beat rate cannot be negative. If \(\omega_-\) and
\(\omega_+\) are almost equal, their difference is small, so the energy exchange
is slow. If the two normal frequencies are farther apart, the beat envelope is
faster.

\begin{quicknote}
    Beats in coupled oscillators and beats in sound have the same mathematics:
    two nearby frequencies combine into rapid oscillation with a slow envelope.
\end{quicknote}

\section{Transient and Steady-State Motion}

A damped, driven oscillator is the standard model for resonance and transients.
It obeys

\begin{namedformula}{Driven damped oscillator}
    m\ddot{x}+b\dot{x}+kx=F_0\cos(\omega_d t)
\end{namedformula}

The solution has two parts:

\[
    x(t)=x_{\text{transient}}(t)+x_{\text{steady}}(t).
\]

The transient part depends strongly on initial conditions and dies away because
of damping. The steady-state part survives after a long time and oscillates at
the driving angular frequency \(\omega_d\), not necessarily at the natural
frequency.

For an underdamped oscillator, the transient part may be written as

\[
    x_{\text{transient}}(t)
    =Ae^{-\gamma t}\cos(\omega' t+\phi),
    \qquad
    \gamma=\frac{b}{2m},
\]

where its damped natural angular frequency is

\[
    \omega'
    =\sqrt{\omega_0^2-\gamma^2}
    =\sqrt{\frac{k}{m}-\left(\frac{b}{2m}\right)^2}.
\]

Here, \(A\) and \(\phi\) are determined by the initial conditions, while the
exponential factor \(e^{-\gamma t}\) makes the transient amplitude decay. The
frequency \(\omega'\) is fixed by the oscillator's mass \(m\), damping constant
\(b\), and spring constant \(k\). It is not generally the same as \(\omega_d\),
which is selected by the external driving source.

Combining the two parts, the complete underdamped motion has the form

\[
    x(t)
    =\underbrace{Ae^{-\gamma t}\cos(\omega' t+\phi)}_
        {\text{transient at the damped natural frequency}}
    +\underbrace{X(\omega_d)\cos(\omega_d t-\delta)}_
        {\text{steady state at the driving frequency}}.
\]

The two frequencies may be close when the oscillator is driven near its damped
natural frequency, but they still describe different physical quantities.

\begin{importantnote}
    Notation varies between textbooks. Some authors use \(\omega_d\) for the
    damped natural frequency. In these notes, \(\omega_d\) always means the
    \textbf{driving angular frequency}, while \(\omega'\) means the
    \textbf{damped natural angular frequency}.
\end{importantnote}

\begin{namedformula}{Steady-state response}
    x_{\text{steady}}(t)=X(\omega_d)\cos(\omega_d t-\delta)
\end{namedformula}

In formula F3.2, \(x_{\text{steady}}(t)\) is the displacement after the
transient part has died away. The symbol \(t\) is time. The factor
\(X(\omega_d)\) is the steady-state amplitude, and it depends on the driving
angular frequency \(\omega_d\). The cosine term shows that the final motion
oscillates at the same angular frequency as the external driving force. The
quantity \(\delta\) is the phase lag: it tells how far behind the driving force
the oscillator's motion is. Thus the formula says that the long-term motion is
sinusoidal, but its size and timing are controlled by the driving frequency,
damping, mass, and spring stiffness.

\begin{namedformula}{Amplitude response}
    X(\omega_d)=
    \frac{F_0}
    {\sqrt{(k-m\omega_d^2)^2+(b\omega_d)^2}}
\end{namedformula}

Formula F3.3 explains the amplitude \(X(\omega_d)\) used in F3.2. The numerator
\(F_0\) is the maximum value of the driving force. A larger \(F_0\) gives a
larger response. The first term in the denominator, \(k-m\omega_d^2\), compares
the spring restoring effect \(k\) with the inertial effect \(m\omega_d^2\). Near
resonance these nearly cancel, so the denominator becomes small. The second
term, \(b\omega_d\), is the damping contribution; it prevents the amplitude from
becoming infinite.

\begin{namedformula}{Phase lag}
    \tan\delta=
    \frac{b\omega_d}{k-m\omega_d^2}
\end{namedformula}

Formula F3.4 gives the phase lag \(\delta\). The numerator \(b\omega_d\) shows
that damping and driving frequency increase the tendency for the motion to lag
behind the force. The denominator \(k-m\omega_d^2\) again compares spring
response with inertia. At low frequency, \(k\) dominates and the lag is small.
Near resonance, \(k-m\omega_d^2\) is close to zero and the lag is about
\(\pi/2\). At very high frequency, inertia dominates and the motion is almost
opposite in phase to the force.

At low driving frequency, the oscillator almost follows the force. Near
resonance, the amplitude becomes large and the phase lag is about \(\pi/2\). At
very high driving frequency, inertia dominates and the oscillator cannot keep up
with the driver.

\begin{definition}
    A \textbf{transient} is the part of the motion that depends on initial
    conditions and fades with time. The \textbf{steady state} is the long-term
    forced motion after transients have decayed.
\end{definition}

\subsection{Resonance}

The undamped natural angular frequency is

\begin{namedformula}{Natural angular frequency}
    \omega_0=\sqrt{\frac{k}{m}}
\end{namedformula}

For light damping, the displacement amplitude is largest near

\begin{namedformula}{Approximate resonance condition}
    \omega_d\approx\omega_0
\end{namedformula}

More carefully, the frequency that maximizes displacement amplitude is slightly
below \(\omega_0\) when damping is present:

\begin{namedformula}{Damped displacement resonance}
    \omega_{\text{res}}
    =
    \sqrt{\omega_0^2-\frac{b^2}{2m^2}}
\end{namedformula}

\begin{importantnote}
    Damping prevents the amplitude from becoming infinite. Without damping, the
    ideal driven oscillator has unbounded growth at exact resonance.
\end{importantnote}

\subsection{Worked Example: Driven Response}

A mass-spring system has \(m=\SI{1.0}{kg}\), \(k=\SI{100}{N.m^{-1}}\), damping
\(b=\SI{2.0}{kg.s^{-1}}\), and driving force amplitude \(F_0=\SI{5.0}{N}\). Find
the response amplitude when \(\omega_d=\SI{10}{rad.s^{-1}}\).

Here

\[
    k-m\omega_d^2=100-1(10)^2=0.
\]

Thus the amplitude is controlled by the damping term:

\[
    X=\frac{5.0}{\sqrt{0^2+(2.0\cdot10)^2}}
      =\frac{5.0}{20}
      =\SI{0.25}{m}.
\]

\[
    \boxed{X=\SI{0.25}{m}}
\]

The phase lag is

\[
    \tan\delta\to\infty,
    \qquad
    \delta\approx\frac{\pi}{2}.
\]

So near resonance the displacement lags the driving force by about a quarter
cycle.

\section{Interference of Waves}

Interference is superposition with phase awareness. When two waves overlap, the
medium does not choose one wave or the other; it adds their displacements.

\begin{definition}
    The \textbf{principle of superposition} says that the net displacement is the
    sum of the displacements that each wave would have produced separately.
\end{definition}

For two sinusoidal waves of the same amplitude, same wavelength, and same
angular frequency moving in the same direction,

\[
    y_1=y_m\sin(kx-\omega t),
    \qquad
    y_2=y_m\sin(kx-\omega t+\phi).
\]

Their sum is another sinusoidal wave:

\begin{namedformula}{Two-wave interference}
    y'=2y_m\cos\left(\frac{\phi}{2}\right)
    \sin\left(kx-\omega t+\frac{\phi}{2}\right)
\end{namedformula}

The resultant amplitude is

\begin{namedformula}{Resultant amplitude}
    y'_m=2y_m\left|\cos\left(\frac{\phi}{2}\right)\right|
\end{namedformula}

\begin{itemize}
    \item If \(\phi=0\), the waves are in phase and interference is fully
    constructive.
    \item If \(\phi=\pi\), the waves are out of phase and interference is fully
    destructive for equal amplitudes.
    \item If \(\phi\) is between these values, the interference is partial.
\end{itemize}

\subsection{Path Difference}

When two coherent waves meet at a point after travelling different distances,
their phase difference is linked to path difference:

\begin{namedformula}{Phase from path difference}
    \phi=\frac{2\pi\Delta L}{\lambda}
\end{namedformula}

\begin{namedformula}{Constructive and destructive path difference}
    \Delta L=m\lambda
    \quad\text{constructive},
    \qquad
    \Delta L=\left(m+\frac{1}{2}\right)\lambda
    \quad\text{destructive}
\end{namedformula}

Here \(m=0,1,2,\ldots\), assuming the sources start in phase.

\subsection{Worked Example: Path Difference}

Two speakers emit the same tone in phase. At a listening point, the path length
difference is \(\Delta L=\SI{1.50}{m}\). If the wavelength is
\(\lambda=\SI{0.50}{m}\), then

\[
    \frac{\Delta L}{\lambda}=\frac{1.50}{0.50}=3.
\]

Since the path difference is an integer number of wavelengths, the interference
is constructive.

\[
    \boxed{\text{Constructive interference}}
\]

\section{Diffraction}

Diffraction is the spreading of waves when they pass through an opening or
around an obstacle. It is strongest when the opening size is comparable to the
wavelength. If the opening is much larger than the wavelength, the wave mostly
continues forward with little spreading.

\begin{quicknote}
    Interference is about adding waves. Diffraction is about how one wavefront
    spreads and interferes with itself after being restricted by an aperture.
\end{quicknote}

\subsection{Single-Slit Diffraction}

For a slit of width \(a\), destructive minima occur at

\begin{namedformula}{Single-slit minima}
    a\sin\theta=m\lambda,
    \qquad
    m=1,2,3,\ldots
\end{namedformula}

Single-slit maxima are less simple than minima. The central maximum occurs at
\(\theta=0\). The side maxima are found exactly from

\begin{namedformula}{Single-slit secondary maxima}
    \tan\beta=\beta,
    \qquad
    \beta=\frac{\pi a\sin\theta}{\lambda}
\end{namedformula}

For many introductory calculations, the side maxima are estimated by placing
them roughly halfway between neighbouring minima:

\begin{namedformula}{Approximate single-slit secondary maxima}
    a\sin\theta\approx\left(m+\frac{1}{2}\right)\lambda,
    \qquad
    m=1,2,3,\ldots
\end{namedformula}

This maxima formula is approximate, while the minima formula above is exact for
the ideal single-slit pattern. The approximation becomes most useful when only
the angular positions of the side bright fringes are needed.

The central maximum is wide when \(a\) is small or \(\lambda\) is large. That is
why long-wavelength waves, such as sound, bend noticeably around doorways, while
visible light usually does not bend dramatically around everyday openings.

\subsection{Circular Aperture and Resolution}

For a circular aperture of diameter \(D\), the first dark ring of the diffraction
pattern occurs at approximately

\begin{namedformula}{Rayleigh angular scale}
    \theta\approx1.22\frac{\lambda}{D}
\end{namedformula}

This sets a diffraction limit: even a perfect lens cannot resolve two point
sources if diffraction spreads their images too much.

\subsection{Worked Example: Single-Slit Minimum}

Light of wavelength \(\SI{600}{nm}\) passes through a slit of width
\(\SI{0.20}{mm}\). Find the first minimum angle.

\[
    a\sin\theta=\lambda
    \quad\Longrightarrow\quad
    \sin\theta=\frac{600\times10^{-9}}{0.20\times10^{-3}}.
\]

\[
    \sin\theta=3.0\times10^{-3},
    \qquad
    \theta\approx3.0\times10^{-3}\text{ rad}.
\]

\[
    \boxed{\theta\approx\SI{0.172}{\degree}}
\]

\section{Dispersion}

Dispersion occurs when wave speed depends on frequency or wavelength. In a
nondispersive medium, all sinusoidal components travel with the same speed, so a
pulse preserves its shape. In a dispersive medium, different components travel at
different speeds, so a pulse can spread and change shape.

\begin{definition}
    A \textbf{dispersion relation} is a formula connecting angular frequency
    \(\omega\) and wave vector magnitude \(k\).
\end{definition}

\begin{namedformula}{Phase and group velocity}
    v_p=\frac{\omega}{k},
    \qquad
    v_g=\frac{d\omega}{dk}
\end{namedformula}

The distinction is easiest to see by imagining a short pulse travelling along a
rope. The pulse can be constructed by adding many sinusoidal waves with nearby
frequencies and wave numbers. Inside the pulse, individual crests and troughs
form a rapidly repeating pattern, while the whole region of noticeable
disturbance forms a slower-changing group.

The phase velocity \(v_p\) is the speed of an individual feature of the repeating
pattern. If one particular crest is followed as it moves through the pulse, its
speed is \(v_p\). The group velocity \(v_g\) is the speed of the pulse as a whole:
it is found by following the region in which the oscillations are concentrated.

An analogy is a group of people walking through a moving spotlight. Each person
represents a crest, while the illuminated group represents the wave packet. A
person can enter the illuminated region at the back, pass through it, and leave
at the front. In the same way, crests can move through a wave packet even when
the packet itself travels at a different speed. Thus, \(v_p\) describes the
motion of the carrier oscillations, whereas \(v_g\) describes the motion of the
overall disturbance or modulation. Group velocity is often the speed at which
wave energy and information are transported.

The formula

\[
    v_p=\frac{\omega}{k}=f\lambda
\]

measures how quickly the repeating phase pattern of one sinusoidal component
moves. The derivative

\[
    v_g=\frac{d\omega}{dk}
\]

measures how the angular frequencies of neighbouring components change as their
wave numbers change. That change determines how quickly their region of
reinforcement moves.

In strongly dispersive or absorbing media, group velocity can behave unusually
and does not always equal the true arrival speed of new information. For the
ordinary weakly dispersive media considered here, however, it is normally the
relevant speed of the pulse and its energy.

\subsection{Calculating Group Velocity from a Dispersion Relation}

The method depends on the form in which the dispersion relation is given. If
angular frequency is explicitly written as a function of wave number,

\[
    \omega=\omega(k),
\]

differentiate it directly:

\begin{namedformula}{Group velocity from \(\omega(k)\)}
    v_g=\frac{d\omega}{dk}
\end{namedformula}

For example, if \(\omega=Ak^2\), then

\[
    v_g=\frac{d}{dk}(Ak^2)=2Ak.
\]

If wave number is instead written as a function of angular frequency,

\[
    k=k(\omega),
\]

calculate \(dk/d\omega\) and take its reciprocal:

\begin{namedformula}{Group velocity from \(k(\omega)\)}
    v_g=\left(\frac{dk}{d\omega}\right)^{-1}
\end{namedformula}

This follows from the inverse-derivative relation

\[
    \frac{d\omega}{dk}\frac{dk}{d\omega}=1.
\]

For example, if \(k=B\omega^2\), then

\[
    \frac{dk}{d\omega}=2B\omega,
    \qquad
    v_g=\frac{1}{2B\omega}.
\]

These are two alternative procedures: either find \(d\omega/dk\) directly, or
find \(dk/d\omega\) and take its reciprocal. No additional derivative is
multiplied after taking the reciprocal.

Sometimes the dispersion relation is implicit, so that \(\omega\) cannot easily
be isolated. If it is written as

\[
    F(\omega,k)=0,
\]

differentiate the whole relation with respect to \(k\):

\[
    \frac{\partial F}{\partial\omega}\frac{d\omega}{dk}
    +\frac{\partial F}{\partial k}=0.
\]

Solving for the required derivative gives

\begin{namedformula}{Group velocity from an implicit relation}
    v_g=\frac{d\omega}{dk}
       =-\frac{\partial F/\partial k}{\partial F/\partial\omega}
\end{namedformula}

Phase velocity does not require differentiation; once \(\omega\) and \(k\) are
known, use \(v_p=\omega/k\). A useful summary is: divide to obtain phase velocity
and differentiate to obtain group velocity.

\begin{importantnote}
    In a nondispersive medium, \(\omega=vk\), so \(v_p=v_g=v\). In a dispersive
    medium, \(v_p\) and \(v_g\) are generally different.
\end{importantnote}

\subsection{Worked Example: A Dispersive Wave}

Suppose a medium has dispersion relation

\[
    \omega=\alpha k^2,
\]

where \(\alpha=\SI{0.040}{m^2.s^{-1}}\). At \(k=\SI{25}{m^{-1}}\),

\[
    v_p=\frac{\omega}{k}=\frac{\alpha k^2}{k}=\alpha k
       =(0.040)(25)=\SI{1.0}{m.s^{-1}}.
\]

\[
    v_g=\frac{d\omega}{dk}=2\alpha k
       =2(0.040)(25)=\SI{2.0}{m.s^{-1}}.
\]

\[
    \boxed{v_p=\SI{1.0}{m.s^{-1}},
    \qquad
    v_g=\SI{2.0}{m.s^{-1}}}
\]

The envelope travels twice as fast as the individual crests.

\section{Polarization}

Polarization describes the direction of oscillation of a transverse wave. It is
meaningful for transverse waves such as electromagnetic waves and waves on a
string. It is not available for an ideal longitudinal sound wave in air, because
the displacement is along the direction of propagation.

For an electromagnetic wave travelling in the \(+z\) direction, the electric
field may be linearly polarized in the \(x\) direction:

\[
    \vec{E}(z,t)=E_0\cos(kz-\omega t)\,\hat{x}.
\]

A polarizer transmits the component of the electric field along its transmission
axis. Since intensity is proportional to field amplitude squared, the transmitted
intensity follows Malus's law.

\begin{namedformula}{Malus's law}
    I=I_0\cos^2\theta
\end{namedformula}

Here \(\theta\) is the angle between the incident polarization direction and the
polarizer axis.

\subsection{Types of Polarization}

Linear polarization has oscillation along a fixed line. Circular polarization can
be formed by two perpendicular components of equal amplitude with a phase
difference of \(\pi/2\). Elliptical polarization is the general case, where the
tip of the field vector traces an ellipse.

\begin{quicknote}
    Polarization tells you about the direction of the disturbance, not the
    direction of wave travel.
\end{quicknote}

\subsection{Worked Example: Two Polarizers}

Unpolarized light of intensity \(I_{\text{in}}\) passes through a polarizer, then
through a second polarizer at \(\theta=60^\circ\) to the first. The first
polarizer transmits half the intensity:

\[
    I_1=\frac{1}{2}I_{\text{in}}.
\]

The second follows Malus's law:

\[
    I_2=I_1\cos^2(60^\circ)
       =
       \frac{1}{2}I_{\text{in}}\left(\frac{1}{2}\right)^2.
\]

\[
    \boxed{I_2=\frac{1}{8}I_{\text{in}}}
\]

\section{Mathematical Modelling with Wave Equations}

The wave equation is a mathematical statement that curvature in space produces
acceleration in time. In one dimension, it is

\begin{namedformula}{One-dimensional wave equation}
    \frac{\partial^2 u}{\partial t^2}
    =
    v^2
    \frac{\partial^2 u}{\partial x^2}
\end{namedformula}

Here \(u(x,t)\) may represent transverse displacement of a string, pressure
variation in a sound wave, or one component of an electromagnetic field. The
quantity \(v\) is the wave speed in the medium.

\subsection{General Solutions in One Dimension}

The one-dimensional wave equation admits right-moving and left-moving solutions:

\begin{namedformula}{D'Alembert form}
    u(x,t)=f(x-vt)+g(x+vt)
\end{namedformula}

The function \(f(x-vt)\) travels in the \(+x\) direction. The function
\(g(x+vt)\) travels in the \(-x\) direction.

\subsection{The Two-Dimensional Wave Equation}

For a membrane or water-surface model with displacement \(u(x,y,t)\),

\begin{namedformula}{Two-dimensional wave equation}
    \frac{\partial^2 u}{\partial t^2}
    =
    v^2
    \left(
    \frac{\partial^2 u}{\partial x^2}
    +
    \frac{\partial^2 u}{\partial y^2}
    \right)
\end{namedformula}

The expression in parentheses is the two-dimensional Laplacian,
\(\nabla^2u\). Boundary conditions now matter strongly. A rectangular membrane,
a circular drumhead, and an open water surface can have very different mode
patterns even if the local wave speed is similar.

\subsection{The Three-Dimensional Wave Equation}

For waves in space,

\begin{namedformula}{Three-dimensional wave equation}
    \frac{\partial^2 u}{\partial t^2}
    =
    v^2\nabla^2u
    =
    v^2
    \left(
    \frac{\partial^2 u}{\partial x^2}
    +
    \frac{\partial^2 u}{\partial y^2}
    +
    \frac{\partial^2 u}{\partial z^2}
    \right)
\end{namedformula}

Plane waves are solutions of the form

\begin{namedformula}{Plane wave in three dimensions}
    u(\vec{r},t)=A\cos(\vec{k}\cdot\vec{r}-\omega t+\phi)
\end{namedformula}

The dispersion relation for a nondispersive three-dimensional wave is

\begin{namedformula}{Nondispersive relation}
    \omega=v|\vec{k}|
\end{namedformula}

\subsection{Initial and Boundary Conditions}

A wave equation alone is not a full physical prediction. To determine the motion,
we also need initial conditions and boundary conditions. Initial conditions say
what the medium is doing at \(t=0\). Boundary conditions say how the edges or
interfaces behave.

\begin{itemize}
    \item Fixed end: displacement is zero at the boundary.
    \item Free end: slope or force condition is zero at the boundary.
    \item Interface: displacement and force, or field components, must match
    according to the medium.
\end{itemize}

\subsection{Worked Example: Testing a Plane Wave}

Show that

\[
    u(x,t)=A\cos(kx-\omega t)
\]

solves the one-dimensional wave equation if \(\omega=vk\).

Differentiate twice:

\[
    \frac{\partial^2u}{\partial t^2}
    =
    -\omega^2A\cos(kx-\omega t),
\]

\[
    \frac{\partial^2u}{\partial x^2}
    =
    -k^2A\cos(kx-\omega t).
\]

The wave equation requires

\[
    -\omega^2A\cos(kx-\omega t)
    =
    v^2[-k^2A\cos(kx-\omega t)].
\]

Cancel the common factor:

\[
    \omega^2=v^2k^2,
    \qquad
    \omega=vk.
\]

\[
    \boxed{\text{The sinusoidal wave solves the equation when } \omega=vk.}
\]

\section{Propagation in Different Media}

The wave equation tells us the mathematical shape of propagation, but the medium
sets the wave speed and often the allowed wave type. Different media resist
deformation in different ways and carry inertia differently.

\subsection{String Waves}

For a stretched string,

\begin{namedformula}{String wave speed}
    v=\sqrt{\frac{\tau}{\mu}}
\end{namedformula}

The tension \(\tau\) supplies the restoring effect and the linear density
\(\mu\) supplies inertia. More tension gives faster waves; more mass per length
gives slower waves.

\subsection{Sound in Fluids}

For sound in a fluid,

\begin{namedformula}{Sound speed in a fluid}
    v=\sqrt{\frac{B}{\rho}}
\end{namedformula}

The bulk modulus \(B\) measures resistance to compression, while \(\rho\) is
density. Stiffer media tend to carry sound faster; denser media tend to carry it
slower when stiffness is fixed.

\subsection{Sound in Solids}

Solids can support longitudinal and transverse mechanical waves. A simplified
longitudinal speed in a slender rod is

\begin{namedformula}{Longitudinal wave speed in a rod}
    v=\sqrt{\frac{Y}{\rho}}
\end{namedformula}

Here \(Y\) is Young's modulus. Real solids can also support shear waves, surface
waves, and multiple polarizations of elastic motion.

\subsection{Electromagnetic Waves}

In vacuum, electromagnetic waves travel at

\begin{namedformula}{Speed of light}
    c=\frac{1}{\sqrt{\mu_0\epsilon_0}}
\end{namedformula}

In a transparent medium, a common model is

\begin{namedformula}{Light speed in a medium}
    v=\frac{c}{n}
\end{namedformula}

where \(n\) is the refractive index. If \(n\) depends on wavelength, the medium
is dispersive and separates colours.

\subsection{Water Waves}

Water waves can be dispersive. In deep water, gravity waves approximately obey

\begin{namedformula}{Deep-water gravity waves}
    \omega^2=gk
\end{namedformula}

Thus

\[
    v_p=\frac{\omega}{k}=\sqrt{\frac{g}{k}},
    \qquad
    v_g=\frac{d\omega}{dk}=\frac{1}{2}\sqrt{\frac{g}{k}}.
\]

The group velocity is half the phase velocity, so wave packets and individual
crests do not move together.

\subsection{Worked Example: Comparing Media}

Compare a string wave and a sound wave. A string has
\(\tau=\SI{100}{N}\) and \(\mu=\SI{0.010}{kg.m^{-1}}\). A fluid has
\(B=\SI{2.2e9}{Pa}\) and \(\rho=\SI{1000}{kg.m^{-3}}\).

For the string,

\[
    v_{\text{string}}=\sqrt{\frac{100}{0.010}}=\SI{100}{m.s^{-1}}.
\]

For the fluid,

\[
    v_{\text{sound}}=\sqrt{\frac{2.2\times10^9}{1000}}
    \approx\SI{1.48e3}{m.s^{-1}}.
\]

\[
    \boxed{v_{\text{string}}=\SI{100}{m.s^{-1}},
    \qquad
    v_{\text{sound}}\approx\SI{1480}{m.s^{-1}}}
\]

The fluid wave is much faster here because the fluid is very resistant to
compression compared with its density.

\section{Connecting the Topics}

Resonance, interference, diffraction, dispersion, and polarization are not
separate islands. They are behaviours of solutions to wave and oscillator
equations.

Resonance happens when a system is driven near one of its natural modes.
Transients describe how the system approaches its long-term motion. Coupling
creates normal modes, and a continuous set of coupled oscillators gives a wave
equation. Interference is the linear addition of solutions. Diffraction is the
spreading forced by apertures and boundaries. Dispersion appears when the medium
changes the relationship between \(\omega\) and \(k\). Polarization records the
directional structure of transverse waves.

\begin{importantnote}
    When solving a problem, first identify the model: oscillator, coupled
    oscillator, travelling wave, standing wave, aperture diffraction, dispersive
    packet, polarized transverse wave, or wave propagation in a specified medium.
    The correct equation usually becomes much clearer after that classification.
\end{importantnote}

\section{Compact Formula Map}

\begin{multicols}{2}

\subsection*{Coupled SHM}
\[
    q_+=x_1+x_2,\quad q_-=x_1-x_2
\]
\[
    \omega_+=\sqrt{k/m},
    \quad
    \omega_-=\sqrt{(k+2k_c)/m}
\]

\subsection*{Driven Damped SHM}
\[
    m\ddot{x}+b\dot{x}+kx=F_0\cos(\omega_d t)
\]
\[
    X=
    \frac{F_0}{\sqrt{(k-m\omega_d^2)^2+(b\omega_d)^2}}
\]
\[
    \omega_0=\sqrt{k/m}
\]

\subsection*{Interference}
\[
    y'_m=2y_m\left|\cos(\phi/2)\right|
\]
\[
    \phi=\frac{2\pi\Delta L}{\lambda}
\]

\subsection*{Diffraction}
\[
    a\sin\theta=m\lambda
\]
\[
    \theta\approx1.22\lambda/D
\]

\subsection*{Dispersion}
\[
    v_p=\omega/k,
    \qquad
    v_g=d\omega/dk
\]

\subsection*{Polarization}
\[
    I=I_0\cos^2\theta
\]

\subsection*{Wave Equations}
\[
    u_{tt}=v^2u_{xx}
\]
\[
    u_{tt}=v^2(u_{xx}+u_{yy})
\]
\[
    u_{tt}=v^2\nabla^2u
\]

\subsection*{Media}
\[
    v=\sqrt{\tau/\mu}
\]
\[
    v=\sqrt{B/\rho}
\]
\[
    v=\sqrt{Y/\rho}
\]
\[
    v=c/n
\]

\end{multicols}

\end{document}
